\documentclass[11pt]{article}
\usepackage[margin=1in]{geometry}
\usepackage{amsmath,amssymb,amsthm,mathtools}
\usepackage{graphicx,microtype,xcolor}
\usepackage[colorlinks=true,linkcolor=blue!55!black,citecolor=blue!55!black,urlcolor=blue!55!black]{hyperref}

\newtheorem{theorem}{Theorem}
\newtheorem{lemma}{Lemma}
\newtheorem{remark}{Remark}
\newcommand{\A}{\mathcal A}
\newcommand{\T}{\mathcal T}
\newcommand{\DI}{C_{\mathrm{DI}}}

\title{An Index Obstruction to Renorming Weighted Bicyclic Algebras}
\author{Automated mathematics run \texttt{fa\_banach\_001}, lane 7}
\date{17 August 2026}

\begin{document}
\maketitle

\begin{center}
\fbox{\parbox{0.92\textwidth}{\textbf{Status.} Full solution, likely valid,
pending human review.  The source's renorming question has a negative answer:
for its weighted bicyclic algebra $\A_x$, every equivalent unital
Banach-algebra norm has $\DI\ge e^{2x}$.  The original weighted norm attains
equality.}}
\end{center}

\section{Source question and result}

Let $BC=\langle p,q:pq=e\rangle$ be the bicyclic monoid.  Every element has a
unique reduced form $q^\alpha p^\beta$, $\alpha,\beta\ge0$.  For $x>0$, Daws
and Horv\'ath define the weight
\[
 \omega_x(s)=
 \begin{cases}
 1,&s=e\text{ or }s=qp,\\
 e^{x(\alpha+\beta)},&s=q^\alpha p^\beta\notin\{e,qp\},
 \end{cases}
\]
and the unital Banach algebra $\A_x=\ell^1(BC,\omega_x)$.  For a unital
Banach algebra $A$ they set
\[
 \DI(A)=\inf\{\|a\|\|b\|:ab=1,\ ba\ne1\}.
\]
They prove that the standard witnesses $\delta_p,\delta_q$ stay large under
every equivalent norm, but leave open whether some other witnesses can make
$\DI(\A_x)$ equal to $1$, or bounded by an absolute constant
\cite[p.~18]{DH2021}:

\begin{figure}[ht]
\centering
\includegraphics[width=\textwidth]{figures/open_question_crop.png}
\caption{The renorming question in \cite[p.~18]{DH2021}.}
\end{figure}

\begin{theorem}[Renorming obstruction]\label{thm:main}
Let $x>0$.  If $\|\cdot\|_0$ is any norm equivalent to the original norm of
$\A_x$ for which $(\A_x,\|\cdot\|_0)$ is a unital Banach algebra with
$\|1\|_0=1$, then
\begin{equation}\label{eq:main}
                  \boxed{\DI(\A_x,\|\cdot\|_0)\ge e^{2x}.}
\end{equation}
For the original weighted norm, equality holds.  Hence there is no absolute
constant $K$ such that every Dedekind-infinite Banach algebra admits an
equivalent unital Banach-algebra norm with $\DI\le K$.
\end{theorem}

The proof combines annular characters with the Fredholm index of a Toeplitz
operator.  The index forces every possible one-sided inverse witness---not
only the displayed generators---to wind nontrivially around the annulus.

\section{Annular characters}

For $z$ in the closed annulus
\[
       \mathbb A_x=\{z\in\mathbb C:e^{-x}\le|z|\le e^x\},
\]
define on the monoid basis
\begin{equation}\label{eq:character}
 \chi_z(\delta_{q^\alpha p^\beta})=z^{\alpha-\beta}.
\end{equation}
The degree $\alpha-\beta$ is additive under multiplication in $BC$, so
$\chi_z$ is multiplicative.  Moreover
\[
 |z^{\alpha-\beta}|
 \le e^{x|\alpha-\beta|}
 \le\omega_x(q^\alpha p^\beta);
\]
the exceptional elements $e,qp$ both have degree zero.  Thus \eqref{eq:character}
extends to a continuous character of $\A_x$.

For $a\in\A_x$ write $\widehat a(z)=\chi_z(a)$.  The defining series converges
uniformly on $\mathbb A_x$ and normally in its interior, so $\widehat a$ is
continuous on the closed annulus and holomorphic inside.  Since equivalent
Banach-algebra norms do not change invertibility or spectra,
\begin{equation}\label{eq:spectral}
 \|a\|_0\ge r_{\A_x}(a)\ge
 \sup_{z\in\mathbb A_x}|\widehat a(z)|.
\end{equation}

\section{Every one-sided inverse has nonzero index}

Let $S$ be the unilateral shift on $\ell^2(\mathbb N_0)$ and define
\begin{equation}\label{eq:toeplitzrep}
 \pi(\delta_q)=S,\qquad \pi(\delta_p)=S^*.
\end{equation}
Because $S^*S=I$, this is a contractive representation of the bicyclic monoid
and extends continuously to $\A_x$.  Its image lies in the Toeplitz algebra
$\T=C^*(S)$, and the Toeplitz symbol of $\pi(a)$ is
\begin{equation}\label{eq:symbol}
             \sigma(\pi(a))(\zeta)=\widehat a(\zeta),
             \qquad |\zeta|=1.
\end{equation}

We record that $\pi$ is faithful.  Indeed, suppose
$a=\sum_{\alpha,\beta}c_{\alpha\beta}
\delta_{q^\alpha p^\beta}$ and $\pi(a)=0$.  Fix an integer
$d=\alpha-\beta$.  Along the $d$th matrix diagonal, the entry beginning at
column $j$ is the absolutely convergent partial sum
\[
 \sum_{\substack{\beta\le j\\\beta+d\ge0}}c_{\beta+d,\beta}.
\]
All these sums vanish; taking successive differences gives
$c_{\beta+d,\beta}=0$ for every admissible $\beta$.  Varying $d$ recovers all
coefficients.

Now take arbitrary $a,b\in\A_x$ with
\begin{equation}\label{eq:onesided}
                         ab=1,\qquad ba\ne1.
\end{equation}
Set $A=\pi(a)$ and $B=\pi(b)$.  Then $AB=I$, while faithfulness gives
$BA\ne I$.  Hence $A$ is surjective but not injective.  On the other hand,
\[
                     \widehat a(z)\widehat b(z)=1
                     \quad(z\in\mathbb A_x),
\]
so the Toeplitz symbol of $A$ is nowhere zero on the unit circle.  Therefore
$A$ is Fredholm, and
\[
 \operatorname{ind}A=\dim\ker A>0.
\]
The Toeplitz index theorem \cite[Chapter 7]{Douglas1972} gives
\begin{equation}\label{eq:index}
 \operatorname{ind}A=-\operatorname{wind}
       (\widehat a|_{\mathbb T},0).
\end{equation}
In particular the integer
$k=\operatorname{wind}(\widehat a|_{\mathbb T},0)$ is nonzero.

\section{Winding forces a large annular condition number}

\begin{lemma}[Annular oscillation]\label{lem:annulus}
Let $f$ be continuous and nonvanishing on $\mathbb A_x$, holomorphic in its
interior, and let its winding number around zero be $k$.  Then
\begin{equation}\label{eq:annulus}
 \left(\sup_{\mathbb A_x}|f|\right)
 \left(\sup_{\mathbb A_x}|f|^{-1}\right)
 \ge e^{2x|k|}.
\end{equation}
\end{lemma}

\begin{proof}
The function $z^{-k}f(z)$ has winding zero and hence admits a holomorphic
logarithm $g$ on the open annulus.  Thus
\[
                 \log|f(re^{it})|=k\log r+\Re g(re^{it}).
\]
The circular mean of $\Re g(re^{it})$ is independent of $r$: this follows
from the Laurent expansion of $g$, whose nonconstant terms have mean zero.
Consequently the circular means of $\log|f|$ on the circles of radii $e^x$
and $e^{-x}$ differ by $2kx$.  If
$M=\max_{\mathbb A_x}\log|f|$ and
$m=\min_{\mathbb A_x}\log|f|$, then
$M-m\ge2x|k|$.  Exponentiation gives \eqref{eq:annulus}.
\end{proof}

Apply Lemma~\ref{lem:annulus} to $f=\widehat a$.  Since
$\widehat b=1/\widehat a$, equations \eqref{eq:spectral} and
\eqref{eq:index} yield
\[
 \|a\|_0\|b\|_0
 \ge\left(\sup_{\mathbb A_x}|\widehat a|\right)
      \left(\sup_{\mathbb A_x}|\widehat b|\right)
 \ge e^{2x|k|}\ge e^{2x}.
\]
This holds for every pair in \eqref{eq:onesided}; taking the infimum proves
\eqref{eq:main}.  In the original weighted norm,
$\delta_p\delta_q=1$ and $\delta_q\delta_p\ne1$, while
\[
 \|\delta_p\|_{\omega_x}=\|\delta_q\|_{\omega_x}=e^x.
\]
Thus $\DI(\A_x,\|\cdot\|_{\omega_x})=e^{2x}$, completing the proof of
Theorem~\ref{thm:main}. \qed

\section{Verification, scope, and novelty boundary}

The proof uses only the algebra and weight defined in the source plus the
standard Toeplitz index theorem.  The two points that close the authors'
stated gap are faithfulness of the shift representation and the fact that a
nonzero index enforces annular, rather than merely unit-circle, oscillation.
The conclusion is invariant under arbitrary equivalent unital
Banach-algebra norms because it is expressed through characters and spectral
radii.

The result answers both formulations in the quoted passage: $\DI$ cannot be
renormed to $1$, and no absolute bound independent of $x$ exists.  It also
determines the exact value in the original norm, improving the source's lower
estimate $(e^x/86)^{1/3}$ for this family.

The run indexes, exact arXiv id and title, and arXiv searches for the
Dedekind-infinite renorming terminology were checked.  No later answer or
duplicate packet was found.  Novelty confidence is moderate: the individual
Toeplitz and annulus ingredients are classical, but their use to control
every $\DI$ witness in this weighted algebra was not located.  Human review
should focus on the faithful extension of \eqref{eq:toeplitzrep} and the
application of the Toeplitz index theorem to arbitrary elements of $\A_x$.

\begin{thebibliography}{9}
\bibitem{DH2021}
M. Daws and B. Horv\'ath,
\emph{Ring-theoretic (in)finiteness in reduced products of Banach algebras},
Canad. J. Math. 73 (2021), 1423--1458; arXiv:1912.07108.

\bibitem{Douglas1972}
R. G. Douglas,
\emph{Banach Algebra Techniques in Operator Theory},
Academic Press, 1972.
\end{thebibliography}

\end{document}
